\section{Necesidad de una soluci\'on tecnol\'ogica para productores de La Libertad}

\subsection{Necesidad de anticipar precios FOB}

Los productores de palta Hass en La Libertad requieren anticipar el precio FOB de exportaci\'on para alinear expectativas de ingreso con decisiones de siembra, manejo de campa\~na y negociaci\'on con compradores \citep{VitteryFlores2023, Velasquez2025}. Sin pron\'osticos a horizontes de 4, 6 y 8 semanas ---compatibles con planificaci\'on operativa--- el productor reacciona ex post a ca\'idas de precio ya materializadas, cuando las acciones correctivas son limitadas.

La predicci\'on FOB constituye el n\'ucleo anal\'itico del sistema propuesto (soluci\'on S-3), mediante la comparaci\'on de SARIMA, LSTM y SVR sobre datos Aduanet. \citet{Kurnadipare2025} demuestran que modelos bien calibrados superan extrapolaciones ingenuas en series de exportaci\'on agr\'icola, aunque la transferencia a palta peruana exige validaci\'on local.

Anticipar precios no elimina la incertidumbre, pero permite cuantificarla: rangos esperados, escenarios adversos y probabilidad impl\'icita de m\'argenes negativos. Esta informaci\'on es especialmente valiosa en a\~nos de sobreoferta global o cuando m\'ultiples regiones peruanas concentran cosecha en las mismas semanas \citep{SierraSelva2020, ComexPeru2025}.

\subsection{Necesidad de evaluar escenarios de rentabilidad}

La rentabilidad de la siembra depende de la interacci\'on entre precio FOB proyectado, rendimiento esperado, costos de producci\'on (agua, fertilizantes, mano de obra, cosecha) y costos log\'isticos \citep{Capunay2025, VargasCordova2025}. \citet{Capunay2025} muestran que per\'iodos de precios aparentemente favorables pueden ocultar m\'argenes estrechos cuando los costos unitarios aumentan m\'as r\'apidamente que los ingresos.

Un simulador de escenarios permite al productor evaluar condiciones optimistas, esperadas y pesimistas variando precio, volumen exportable y costos. La soluci\'on S-2 del proyecto traduce predicciones en m\'argenes proyectados y umbrales de viabilidad econ\'omica, cerrando la brecha entre pron\'ostico t\'ecnico y decisi\'on financiera.

\citet{VargasCordova2025} evidencia que productores sin herramientas de simulaci\'on dependen de estimaciones mentales sesgadas, sobreestimando ingresos en a\~nos buenos y subestimando resiliencia necesaria en a\~nos malos. La Libertad, por su concentraci\'on productiva, amplifica externalidades de decisiones colectivas no coordinadas: muchos productores expandiendo siembra simult\'aneamente puede deprimir precios futuros para todos.

\subsection{Importancia de disponer de recomendaciones de siembra}

M\'as all\'a de gr\'aficos y m\'etricas, el productor necesita recomendaciones accionables: ``siembra viable'', ``postergar expansi\'on'', ``reducir inversi\'on en insumos este ciclo''. Estas recomendaciones deben derivarse de reglas transparentes vinculadas a escenarios simulados, no de cajas negras incomprensibles.

\citet{Theofilou2025} se\~nala que la utilidad percibida de modelos predictivos agr\'icolas aumenta cuando se combinan con interfaces comprensibles y contexto local. \citet{ReisFilho2022} enfatizan enriquecer series con informaci\'on textual y explicaciones, enfoque compatible con reportes de ciclos (soluci\'on S-5) que identifiquen periodos hist\'oricamente favorables o adversos para siembra.

Para La Libertad, disponer de recomendaciones basadas en datos Aduanet e INEI democratiza acceso a inteligencia de mercado antes reservada a exportadores con equipos anal\'iticos. El impacto esperado es mejor alineaci\'on entre capacidad productiva regional y se\~nales de demanda global, reduciendo prociclicidad en decisiones de inversi\'on agr\'icola.
